\chapter*{\chapterstyle{V --- Fonctions Vectorielles}} % 95 Fini
\addcontentsline{toc}{section}{Fonctions Vectorielles}

Dans ce chapitre nous étudions les propriétés analytiques et de régularité d'un ensemble plus général de fonctions, les \textbf{fonctions d'une variable à valeurs vectorielles} et plus précisément, on se concentrera principalement sur le cas des fonctions de \(\R\) dans \(\R^n\) muni de sa structure euclidienne si nécessaire.

\subsection*{\subsecstyle{Définition{:}}}
Soit \(A \subset \R\) une partie de \(\R\) et \((f_1, f_2, \ldots, f_n)  \in \mathcal{F}(A, \R)\) des fonctions réelles, alors on appelle \textbf{fonction vectorielle} les fonctions de la forme:
\[
   \begin{aligned}
      f: A &\longrightarrow \R^n\\
      t &\longmapsto (f_1(t), f_2(t), \ldots, f_n(t))
   \end{aligned}
\]
Dans le cas ou \(n = 2\) et si \(\)

\subsection*{\subsecstyle{Limites{:}}}
On définit alors \textbf{la limite d'une fonction vectorielle} \(f\) en un point \(a \in A\), par la limite par composante, ie on a:
\[
   \lim_{x \rightarrow a} f(t) = l \Longleftrightarrow \forall i \in \inticc{1}{n} \; ; \; \lim_{x \rightarrow a} f_i(t) = l_i  
\]
En particulier on remarque alors la simplicité de cette généralisation, les limites se calculent simplement composantes par composantes\footnote[1]{\label{carac}Et on peut alors montrer toutes les caractérisations classiques, séquentielles, topologiques etc..}. On verra par la suite que cette définition permet bien d'étendre toutes les notions analytiques classiques.

\subsection*{\subsecstyle{Propriétés de régularité{:}}}
La limite étant définie composantes par composantes, la notion de continuité s'étend naturellement composantes par composantes et on dira alors (comme dans le cas réel) que \(f\) est continue\footref{carac} si et seulement si:
\[
   \lim_{x \rightarrow a} f(t) = f(a)
\]

Par ailleurs, on étends aussi la notion de dérivabilité et on dit qu'une fonction vectorielle \(f\) est \textbf{dérivable en a} si et seulement si la limite suivante existe:
\[
   \lim_{x \rightarrow a} \frac{f(x) - f(a)}{x - a}
\]
Ce qui équivaut alors à l'existence des limites correspondantes pour les composantes, et \(f\) est donc dérivable en \(a\) (par extension de classe \(\mathcal{C}^k\)) si chacune de ses composantes sont dérivables en \(a\) (ou de classe \(\mathcal{C}^k\) en \(a\)) et on note alors \(f'\) sa \textbf{dérivée} qui s'obtient simplement en dérivant composante par composante.\<

Géométriquement \(f'(t)\) représente \textbf{le vecteur tangent} (aussi appellé vecteur vitesse) à la courbe au temps \(t\). Dans \(\R^2\) par exemple:
\begin{center}
   \begin{tikzpicture}[domain=-2:2, xscale=1.5]
      \tkzInit[xmin=-2,xmax=2,ymin=-2,ymax=2]
      \draw[-latex] (-2,0) -- (2,0) node [right] {$x$};
      \draw[-latex] (0, -2) -- (0,2) node [above] {$y$};
      \tkzFctPar[samples=400,domain=0:2*pi]{1.75*sin(2*t)}{1.75*sin(3*t)};
      \draw[-Stealth, color=BrightRed1, line width=0.35mm] (1.75, 1.23) -- (1.75,2) node [right] {$\;\gamma'(\frac{\pi}{4})$};
   \end{tikzpicture}
\end{center}
\pagebreak
\subsection*{\subsecstyle{Propriétés non-conservées{:}}}
Néanmoins pour l cas de fonctions vectorielles, certaines propriétés ne sont pas conservées. En particulier, \textbf{le théorème de Rolle} n'est plus valable ainsi que le \textbf{le théorème des accroissements finis}.\footnote[1]{Evident aprés un dessin.}
Néanmoins, l'inégalité des accroissement finis reste valide et s'intérprète facilement cinématiquement:
\begin{center}
   \textit{Un point mobile qui se déplace à une vitesse instantanée inférieure à \(K\) sur un temps \(T\), se trouve au maximum à une distance \(KT\) de son point de départ.}
\end{center}

\subsection*{\subsecstyle{Dérivées de composées {:}}}
Soit \(f\) une fonction vectorielle dérivable, on peut alors se demander comment, étant donnée une application linéaire \(L\), se comporte la dérivée de la composée \(L \circ f\) et on montre alors à partir de la définition que:
\[
   (L(f))' = L(f')
\] 
On veut alors généraliser, et on considère deux fonctions vectorielles \(f, g\) dérivables, et une application bilinéaire \(B\), alors de la meme manière on peut montrer que:
\[
   B(f, g)' = B(f', g) + B(f, g')
\]
On peut alors y voir un analogue de la \textbf{règle du produit} pour la dérivation classique, en particulier si \(B = \dotproduct{\cdot}{\cdot}\), on a:
\[
   \dotproduct{f}{g}' = \dotproduct{f'}{g} + \dotproduct{f}{g'}
\]
Finalement, si on a une famille \(f_1, \ldots, f_n\) de fonctions vectorielles dérivables et une application multilinéaire \(M\), alors:
\[
   M(f_1, \ldots, f_n)' = M(f_1', \ldots, f_n) + \ldots + M(f_1, \ldots, f_n')
\]
En particulier si \(M = \det\), et pour les fonctions \(f, g, h\) on a:
\[
   \det(f, g, h)' = \det(f', g, h) + \det(f, g', h) + \det(f, g, h')  
\]

\subsection*{\subsecstyle{Intégration{:}}}
On peut alors raisonnablement penser que l'intégration sur un segment \(\icc{a}{b}\) se généralisera de la meme manière, et c'est le cas ! Si toutes les fonctions \((f_i)\) sont Riemann-intégrables sur le segment, alors on dira que la fonction \(f\) est Riemann-intégrable et on définit:
\[
   \int_{a}^{b} f(t) \d t = \left(\int_{a}^{b} f_1(t) \d t, \ldots, \int_{a}^{b} f_n(t) \d t\right)   
\]
En particulier cela donne lieu à une interprétation cinématique :
\begin{center}
   \textit{Si \(f\) est une fonction qui donne le vecteur vitesse d'un point mobile à un temps \(t\), alors son intégrale est le vecteur position associé.}
\end{center}

\subsection*{\subsecstyle{Points remarquables{:}}}
On peut par la suite définir plus d'outils analytiques sur cet ensemble de fonctions, mais pour cela il faut définir la notion d'arcs paramétré qui nécessite les définitions suivantes qui se ne sont que terminologie:
\begin{itemize}
   \item On appelle \textbf{point régulier} un point \(f(t)\) tel que \(f'(t) \neq 0\)
   \item On appelle \textbf{point critique} un point \(f(t)\) tel que \(f'(t) = 0\)
\end{itemize}
On appelle \textbf{point multiple} un point \(P\) de \(E\) tel qu'il existe \((t_0, \ldots, t_k)\) deux à deux distincts tels que \(\forall i \in \inticc{1}{k} \; f(t_i) = P\) et on appelle alors \textbf{multiplicité} l'entier \(k\).

\chapter*{\chapterstyle{V --- Arcs Paramétrés}} % 95 Fini
\addcontentsline{toc}{section}{Arcs Paramétrés}

Dans ce chapitre, on s'intéresse à un cas particulier de fonctions vectorilles qu'on appelera \textbf{arcs paramétrés}, qui sont des fonctions vectorielles à valeurs dans \(\R^2\) ou \(\R^3\).

\subsection*{\subsecstyle{Définition{:}}}
Soit \(I \subset \R\) un intervalle et \(x, y \in \mathcal{C}^1(I, \R)\), alors on définit la fonction vectorielle suivante:
\[
   \begin{aligned}
      f: I &\longrightarrow \R^2\\
      t &\longmapsto (x(t), y(t))
   \end{aligned}
\]
On apelle alors \textbf{arc paramétré} le couple \((I, f)\) et \textbf{courbe} l'ensemble des points \(f(I)\). Réciproquement, si on a une courbe \(C \subset \R^2\) et qu'il existe une fonction \(f\) telle que \(C = f(I)\), alors \((I, f)\) est apellé \textbf{paramétrage de la courbe}.\<

\subsection*{\subsecstyle{Propriétés{:}}}
Si \(I\) est un intervalle fermé de la forme \(\icc{a}{b}\), alors on appelle \textbf{point initial} (respectivement \textbf{point final}) le point \(f(a)\) (respectivement \(f(b)\)).\<

Pour une courbe donnée, il existe une infinité d'arc paramétrés décrivant cette courbe, par exemple les deux arcs paramétrés suivants sont bien différents mais ont \textbf{même courbe} (le cercle unité).
\begin{align*}
   \lambda_1 &:= \bigl(\icc{0}{2\pi}, (\cos(t), \sin(t))\bigl) \\
   \lambda_2 &:= \bigl(\icc{2\pi}{4\pi}, (\cos(t), \sin(t))\bigl)
\end{align*}
On peut aussi remarquer que tout graphe d'une fonction \(g\) peut être considéré comme la courbe d'un arc paramétré \((I, f)\) où \(f(t) = (t, g(t))\). \<

Si \(f\) est injective, on dit alors que l'arc paramétré est \textbf{simple}\footnote[1]{Qui consiste à dire que la courbe n'a pas de point multiples}. \+
Si \(f'\) ne s'annule pas, on dit alors que l'arc paramétré est \textbf{régulier}\footnote[2]{Qui consiste à dire que la courbe n'a pas de point stationnaires, ie de points ou l'on ``arrête de se déplacer'' sur la courbe.}. \+

On appelle \textbf{courbe fermée} la trajectoire d'une arc paramétré par une fonction périodique.

\subsection*{\subsecstyle{Courbes polaires et paramétrage par l'angle polaire{:}}}
On fixe une origine \(O\) du plan et une demi-droite partant de l'origine.
Tout point du plan peut alors être identifié par ses \textbf{coordonées polaires} \((r, \theta)\) avec \(r\) la distance dirigée\footnote[3]{ On autorise alors une rayon négatif qui signifie simplement que le point a pour coordonées \((r, \pi + \theta)\), ie le point se situe toujours sur la droite d'angle \(\theta\) mais "de l'autre coté" de l'orgine.} entre \(P\) et l'origine et \(\theta\) l'angle entre \(\vv{OP}\) et la demi-droite.\<

Une façon fréquente de définir une courbe est alors de donner son \textbf{equation polaire}\footnote[4]{C'est un cas particulier d'arc paramétré de la forme \((I, (f(\theta)\cos(\theta), f(\theta)\sin(\theta)))\) par passage en coordonées cartésiennes.} de la forme \(r = f(\theta)\) pour \(\theta \in I\). 
\pagebreak

\subsection*{\subsecstyle{Courbes planes remarquables{:}}}
On peut considèrer alors quelques courbes remarquables:
\begin{figure}[h]
   \centering
   \begin{subfigure}{.3\textwidth}
      \centering
      \begin{tikzpicture}[line cap=round]
         \draw[-latex] (0,0) -- (3.5,0) node [right] {$x$};
         \draw[-latex] (0, 0)     -- (0,1) node [above] {$y$};
         \draw[BrightBlue1, thick] plot[variable=\t,domain=0:4*pi,smooth,thick] ({(\t - sin(\t r))/4},{(1 - cos(\t r))/4});
      \end{tikzpicture}
      \caption*{La cycloïde}
   \end{subfigure}\quad
   \begin{subfigure}{.3\textwidth}
      \centering
      \begin{tikzpicture}[line cap=round]
         \draw[-latex] (-1.2,0) -- (1.2,0) node [right] {$x$};
         \draw[-latex] (0,-1)     -- (0,1) node [above] {$y$};
         \draw[BrightBlue1, thick] plot[variable=\t,domain=0:2*pi,smooth,thick] ({(cos(\t r)+1)*(cos(\t r)/2},{(cos(\t r)+1)*(sin(\t r)/2});
      \end{tikzpicture}
      \caption*{Une cardioïde}
   \end{subfigure}\quad
   \begin{subfigure}{.3\textwidth}
      \centering
      \begin{tikzpicture}[line cap=round]
         \draw[-latex] (-1.2,0) -- (1.2,0) node [right] {$x$};
         \draw[-latex] (0,-1.2)     -- (0,1.2) node [above] {$y$};
         \draw[BrightBlue1, thick] plot[variable=\t,domain=0:2*pi,smooth,thick] ({(cos(\t r)^3)},{(sin(\t r)^3)});
      \end{tikzpicture}
      \caption*{Une hypocycloïde}
   \end{subfigure}\quad
   \begin{subfigure}{.3\textwidth}
      \centering
      \begin{tikzpicture}[line cap=round]
         \draw[-latex] (-0.2,0) -- (1.2,0) node [right] {$x$};
         \draw[-latex] (0,-0.75)     -- (0,0.75) node [above] {$y$};
         \draw[BrightBlue1, thick] plot[variable=\t,domain=0:2*pi,smooth,thick] ({(1/1.8+cos(\t r))/2*cos(\t r)},{(1/2+cos(\t r))/1.8*sin(\t r)});
      \end{tikzpicture}
      \caption*{Un limaçon}
   \end{subfigure}\quad
   \begin{subfigure}{.3\textwidth}
      \centering
      \begin{tikzpicture}[line cap=round]
         \draw[-latex] (-1.2,0) -- (1.2,0) node [right] {$x$};
         \draw[-latex] (0,-0.75)     -- (0,0.75) node [above] {$y$};
         \draw[BrightBlue1, thick] plot[variable=\t,domain=0:2*pi,smooth,thick] ({(sqrt(2)/1.5*sin(\t r))/(1+cos(\t r)^2)},{(sqrt(2)/1.5*sin(\t r)*cos(\t r))/(1+cos(\t r)^2)});
      \end{tikzpicture}
      \caption*{Le lemniscate de Bernoulli}
   \end{subfigure}\quad
   \begin{subfigure}{.3\textwidth}
      \centering
      \begin{tikzpicture}[line cap=round]
         \draw[-latex] (-1.2,0) -- (1.2,0) node [right] {$x$};
         \draw[-latex] (0,-0.75)     -- (0,0.75) node [above] {$y$};
         \draw[BrightBlue1, thick] plot[variable=\t,domain=0:2*pi,smooth,thick] ({\t/7*cos(\t r)},{\t/7*sin(\t r)});
      \end{tikzpicture}
      \caption*{La spirale d'Archimède}
   \end{subfigure}\quad
\end{figure}

Ces courbes ont des propriétés trés intéressantes et se retrouvent, par exemple, en physique. En effet, la cycloïde (inversée) représente la pente qui minimise le temps de trajet d'une bille lancée en haut de celle ci. On dit que c'est une courbe \textbf{brachistochrone}.

\subsection*{\subsecstyle{Longueur de courbes paramétrées{:}}}
Soit \((I, f)\) un arc paramétré \textbf{régulier} de classe \(\mathcal{C}^1\). On s'intéresse à la \textbf{longueur de l'arc paramétré} qu'on notera \(L\).\<

On peut alors subdiviser \(I\) en différents intervalles et approximer la longueur de l'arc par la somme des cordes d'extérmités les différentes subdivisions. En faisant alors tendre la longueur des subdivisions vers 0, on obtient, à la limite, une intégrale de Riemann\footnote[1]{L'intégrande est bien intégrable (et même continue) en tant que somme produit et composées de fonctions intégrables (continues)}, et la définition suivante:
\customBox{width=3.5cm}{
   \[
      L := \int_I \vectNorm{f'(t)}\, dt  
   \]
}
On appelle alors \textbf{abcisse curviligne} d'origine \(t_0\) la primitive\footnote[2]{En particulier, on voit que l'abcisse curviligne évaluée en en réel \(t\), est exactement \textbf{la longueur algébrique de l'arc} entre \(t_0\) et \(t\), ce qui explique notamment la dénomination anglophone \textit{arc-length function}.}:
\[
   s(t) = \int_{t_0}^{t} \vectNorm{f'(x)} \d x
\]
Par exemple dans \(\R^2\) pour l'arc paramétré \((I, f)\) avec \(I = \icc{0}{1}\) et \(f(t) = (t, t^2)\), on obtient:
\[
   L := \int_{I} \vectNorm{f'(t)}\,dt = \int_{\icc{0}{1}} \vectNorm{(1, 2t)}\,dt = \int_{\icc{0}{1}}\sqrt{1 + 4t^2}\,dt \approx 1.47
\]
Quand le support de l'arc paramétré est une partie de \(\R^2\), on peut aussi préfèrer exprimer l'arc en coordonées polaires, alors on peut déduire de la définition ci-dessus la caractérisation suivante:
\customBox{width=5cm}{
   \[
      L := \int_I \sqrt{r'(\theta)^2 + r(\theta)^2}\, d\theta 
   \]
}

\pagebreak
\subsection*{\subsecstyle{Aires de courbes paramétrées{:}}}
Soit \((I, f)\) un arc paramétré \textbf{régulier} de classe \(\mathcal{C}^1\) qui a pour courbe une partie de \(\R^2\) et \(I\) un segment. On s'intéresse à \textbf{l'aire entre l'arc paramétré} et l'axe des abcisses qu'on notera \(A\).

On peut alors subdiviser \(I\) et approximer l'aire sous la courbe par la somme d'aires de rectangles. En faisant alors tendre la longueur des subdivisions vers 0, on obtient une intégrale de Riemann\footnote[1]{\label{note1}L'intégrande est bien intégrable en tant que somme produit et composées de fonctions intégrables.}, et la définition suivante:
\customBox{width=3.8cm}{
   \[
      A := \int_I y(t)x'(t)\, dt
   \]
}
Par exemple pour l'arc paramétré \((I, f)\) avec \(I = \icc{0}{1}\) et \(f(t) = (t, t^2)\), on obtient:
\[
   A := \int_I y(t)x'(t)\, dt = \int_{\icc{0}{1}} t^2\, dt = \frac{1}{3}
\]

Dans le cas des courbes en coordonées polaires, on peut calculer l'aire bornée par la courbe, on subdivise alors \(I\) en \((\theta_i)_{i \in \N}\), et on calcule la somme de \textbf{secteurs circulaires} de rayon \(r(\theta_k)\). Or on sait que l'aire d'une cercle de rayon \(r\) est donnée par \(\pi r^2\) donc l'aire d'un secteur circulaire d'angle \(\theta\) est donné par \(A = \frac{1}{2}\theta r^2\).\<

Pour obtenir l'aire délimitée par la courbe polaire on peut alors sommer les aires des secteurs circulaires, on obtient alors, à la limite, une intégrale de Riemann\footnoteref{note1} et la définition suivante:
\customBox{width=4cm}{
   \[
      A := \frac{1}{2} \int_I r(\theta)^2\, d\theta
   \]
}
\begin{center}
   \textit{
      On a alors deux types de calculs d'aires, l'aire des rectangles entre une courbe et l'axe des abcisses et l'aire des secteurs circulaires entre une courbe et l'origne, qui ne sont \textbf{pas du tout équivalentes}.
   }
\end{center}

\subsection*{\subsecstyle{Reparamétrage{:}}}
Soit \(\gamma = (I, f)\) un arc paramétré de classe \(\mathcal{C}^k\). Soit \(\phi\) un \(\mathcal{C}^k\)-diféomorphisme d'un certain intervalle \(J\) dans \(I\), alors on appelle la fonction \(f \circ \phi\) un \textbf{paramétrage admissible} de l'arc et on dit alors que les arcs sont \(\mathcal{C}^k\)\textbf{-équivalents}.\<

Deux arcs \(\mathcal{C}^k\)\textbf{-équivalents} ont même courbes, mêmes tangentes, même longueur d'arcs (entre deux points d'une part, et deux image de l'autres) et même points réguliers\footnote[2]{Ceci justifie de contraindre \(\phi\) à être un difféomorphisme.}.

\subsection*{\subsecstyle{Paramétrage normal{:}}}
Si \(s\) est l'abcisse curviligne d'un arc régulier \(\gamma = (I, f)\) de classe \(\mathcal{C}^1\), alors on peut montrer\footnote[3]{Elle est continue et strictement croissante et sa dérivée ne s'annule pas, donc bijective à réciproque dérivable. Voir le chapitre sur la dérivation, partie ``Dérivée d'une réciproque''.} que \(s\) est un \(\mathcal{C}^1\)-difféomorphisme de \(I\) sur \(J = f(I)\), et en particulier, on peut alors reparamétrer par l'abcisse curviligne:
\customBox{width=3.5cm}{
   \(
      \gamma' = (J, f \circ s^{-1})
   \)
}
On peut alors remarquer alors que ce nouvel arc paramétré s'affranchit de la vitesse de parcours, et parcours la courbe \textbf{à vitesse constante}. Ce reparamétrage appelée \textbf{paramétrage normal} est généralement bien plus naturel quand on s'intéresse à des questions géométriques intrinsèques à la courbe.
